% !TeX root = ../paper.tex

\section{Background and Related Work}
\label{sec:related}

% NOTE (kept deliberately short): the supervisor asked for less weight on the
% background and more on the own contribution. Cite each source at its first
% mention. Replace every % TODO with your own synthesis before submission.

\subsection{Software Quality and ISO/IEC 25010}
% NOTE: only the free ISO/IEC 25010:2023 preview (Scope + functional-suitability
% definitions) is accessible --- the full text is paywalled. Cite iso25010 ONLY
% for the high-level model framing below (all verifiable from the preview Scope).
% The maintainability definition/sub-dimensions are NOT verifiable from ISO here,
% so they are sourced to Heitlager et al. (peer-reviewed, freely available)
% and/or an SE textbook instead.
This work frames software quality per ISO/IEC~25010, which models product quality
as nine characteristics used to specify, measure, and evaluate quality, with
functional correctness being only one subcharacteristic~\citep{iso25010}.
Maintainability is the characteristic of interest here: functional correctness is
necessary but does not, on its own, guarantee that code is easy to understand and
change. Heitlager et al.~\citep{heitlager2007maintainability} argue that
maintainability should be estimated from measures computed directly on the source
code, and document a series of limitations of the Maintainability Index in that
role. This study shares their emphasis on directly measurable source-code
properties, but operationalizes maintainability through the Maintainability Index,
cyclomatic complexity, and cognitive complexity rather than through their SIG
model. The Index is retained deliberately despite their critique: how it behaves
under a feedback loop of this kind is itself one of the questions this study
answers (\cref{sec:results}).

\subsection{Maintainability and Complexity Metrics}
Three structural metrics form the primary measurement basis of this study.
The \emph{Maintainability Index} (MI) combines several base measures---Halstead
volume, cyclomatic complexity, lines of code, and comment density---into a
single composite score~\citep{coleman1994maintainability}. The implementation
used here, \texttt{radon}~6.0.1, applies a normalized 0--100 derivative of that
formula~\citep{radon_docs}. Because the Index is a
composite, a change in it can originate in any of
these inputs, which limits its value as a standalone signal and is why the
underlying components are also reported separately in this study
(\cref{sec:supp}). \emph{Cyclomatic complexity} (CC) counts the
independent paths through a unit of code~\citep{mccabe1976complexity}.
\emph{Cognitive complexity} is a separate metric introduced to address CC's
shortcomings as a measure of understandability: it adds an increment wherever the
code's linear reading order is interrupted, and a further increment for each
level at which such a structure is nested~\citep{cognitivecomplexity}.

This same combination---the Maintainability Index, cyclomatic complexity, cognitive
complexity, Halstead metrics, and lines of code, produced by \texttt{radon} and
\texttt{complexipy}~\citep{complexipy_tool}---has been applied in recent empirical work comparing the code
quality of LLM-generated and human-written Python
programs~\citep{jamil2025llmquality}, providing a relevant empirical precedent for
its use here. A recent multi-dimensional
study of five commercial LLM code-generation tools across five programming languages
shows how such metrics are applied in practice: it
scores generated code on functional correctness, reliability, maintainability, and
complexity, and reports that code smells are pervasive and that some outputs reach
excessive cyclomatic complexity~\citep{li2026multidim}. There, as in the work
discussed below, the metrics assess
code after it has been generated rather than guiding the generation itself.

To capture additional structural and size effects, three supplementary measures
are reported alongside the primary trio: \emph{Halstead volume/effort},
\emph{lines of code} (LOC), and \emph{linter issue counts} (Pylint
refactor/warning/convention)~\citep{pylint_tool}. These are produced by the same tooling at no
additional measurement cost (see \cref{sec:method}). LOC is reported in
particular because the MI incorporates code volume through two separate volume
measures, Halstead volume and lines of code, an ingredient Heitlager et al.\
criticize as lacking any logical
justification~\citep{heitlager2007maintainability}; reporting LOC separately lets
a reader confirm that MI changes are not merely an artifact of code length.

\subsection{LLM Code Generation and Benchmarks}
The MBPP benchmark provides short, self-contained Python programming tasks with
associated tests~\citep{austin2021mbpp}. Together with HumanEval it is, as Blyth
et al.\ observe, a standard benchmark for assessing LLM-generated
code~\citep{blyth2025scam}. This study draws its tasks from the sanitized
subset of MBPP and scores them with MBPP+, the test-augmented variant released
with the EvalPlus tool~\citep{evalplus_tool} (\cref{sec:taskset}).

\subsection{Research Gap}
Blyth et al.~\citep{blyth2025scam} demonstrate that iterative static-analysis
feedback, using Bandit and Pylint reports as prompts, substantially reduces
security, readability, and reliability violations in GPT-4o-generated Python code.
Their loop includes a fitness-based correctness gate ($f(S)=-\infty$ for
test-failing snippets), so passing code cannot be replaced by failing code, and they
report an increase of roughly ten percentage points in the proportion of
functionally correct snippets across the iterations. Their
evaluation measures quality through severity-weighted linter issue counts,
however; structural complexity metrics such as the Maintainability Index,
cyclomatic complexity, and cognitive complexity appear nowhere in their evaluation. A
complementary line of work uses execution results and unit-test outcomes as the
feedback signal~\citep{chen2024selfdebug}; structural quality metrics are not part
of that loop. A further line makes
generation itself quality-aware: ARCHCODE conditions code generation on inferred
functional and non-functional requirements---the latter including maintainability,
robustness, and time and space performance---and ranks candidate solutions by their
compliance with those requirements~\citep{han2024archcode}. This intervenes at
generation time through requirement conditioning rather than by feeding measured
structural metrics back into an iterative loop. ARCHCODE does operationalize
maintainability as a cyclomatic-complexity threshold, checked with the same
\texttt{radon} library used here, but it applies that threshold as a pass/fail
filter over candidate solutions rather than reporting continuous complexity
values as evaluation outcomes, and it considers neither the Maintainability Index
nor cognitive complexity~\citep{han2024archcode}. The study presented
here examines the question these lines leave open: does static-analysis feedback
produce measurable changes in structural complexity metrics on a general
programming benchmark under a frozen, reproducible protocol?
